% =====================================================================
% CAPÍTULO 3: MISIÓN, VISIÓN, VALORES Y OBJETIVOS
% =====================================================================
\chapter{Misión, Visión, Valores y Objetivos}
\label{ch:mision-vision}

\section{Identidad del Sistema}

\begin{cajaNota}[Identidad]
\textbf{Denominación:} Sistema de Gestión Departamental de Informática y
Computación (\sistema).\\[4pt]
\textbf{Naturaleza:} Plataforma integrada de gestión académica, comunicación
institucional y analítica de datos.\\[4pt]
\textbf{Ámbito de aplicación:} \ac{DIC}.\\[4pt]
\textbf{Usuarios:} estudiantes, docentes, secretaría y dirección del departamento.
\end{cajaNota}

\section{Misión}

\begin{cajaObjetivo}[Misión]
Facilitar al \ac{DIC} una plataforma unificada, transparente y trazable que
centralice la gestión de cupos, materias, comunicación, faltas, evaluación
docente, quejas y sugerencias, incorporando analítica de datos que permita
detectar patrones, anticipar problemas y generar soluciones rápidas basadas en
evidencia, con el fin de mejorar la oportunidad de las respuestas, la equidad en
las decisiones y la calidad de la experiencia académica de estudiantes, docentes
y personal administrativo.
\end{cajaObjetivo}

\subsection{Desglose de la misión}
\begin{table}[H]
\centering
\caption{Componentes de la misión}
\label{tab:mision-desglose}
\begin{tabularx}{\textwidth}{@{}l X@{}}
\toprule
\rowcolor{verdeClaro}
\textbf{Componente} & \textbf{Contenido} \\
\midrule
¿Qué hacemos? & Centralizamos y automatizamos los procesos académicos y
administrativos del \ac{DIC}. \\
¿Para quién? & Estudiantes, docentes, secretaría y dirección del departamento. \\
¿Para qué? & Para mejorar oportunidad, trazabilidad, equidad y calidad de las
decisiones académicas. \\
¿Cómo? & Mediante una plataforma integrada con analítica de datos y un motor de
soluciones rápidas. \\
¿Con qué valores? & Transparencia, equidad, confidencialidad, eficiencia,
innovación y servicio. \\
\bottomrule
\end{tabularx}
\end{table}

\section{Visión}

\begin{cajaNota}[Visión]
Ser, en un horizonte de tres años, el sistema de gestión departamental de
referencia en la universidad, reconocido por eliminar las barreras de
comunicación entre estudiantes y docentes, por garantizar evaluaciones docentes
oportunas y por transformar los datos operativos en decisiones académicas
informadas, anticipadas y verificables.
\end{cajaNota}

\subsection{Metas de la visión}
\begin{table}[H]
\centering
\caption{Metas cuantificables asociadas a la visión}
\label{tab:metas-vision}
\small
\begin{tabularx}{\textwidth}{@{}X c c@{}}
\toprule
\rowcolor{azulClaro}
\textbf{Meta} & \textbf{Línea base} & \textbf{Objetivo a 3 años} \\
\midrule
Tiempo de respuesta a solicitudes de cupo & 5--10 días hábiles & $< 24$ horas \\
Evaluaciones docentes completadas en el periodo & $< 40\%$ & $= 100\%$ \\
Quejas con responsable asignado y fecha de cierre & $< 30\%$ & $> 95\%$ \\
Tiempo medio de resolución de queja crítica & $> 15$ días & $< 3$ días \\
Cobertura de la encuesta de horarios & No existente & $> 80\%$ de matriculados \\
Módulos con tablero de verificación & 0 & 100\% de procesos críticos \\
Materias con horario optimizado por modelo & 0 & $> 70\%$ \\
\bottomrule
\end{tabularx}
\end{table}

\section{Valores Institucionales}

\begin{table}[H]
\centering
\caption{Valores y su manifestación en el sistema}
\label{tab:valores}
\small
\begin{tabularx}{\textwidth}{@{}l X@{}}
\toprule
\rowcolor{moradoClaro}
\textbf{Valor} & \textbf{Cómo se manifiesta en el \ac{SGDIC}} \\
\midrule
Transparencia & Estados, plazos y criterios visibles para el solicitante. \\
Equidad & Reglas de priorización publicadas y aplicadas uniformemente. \\
Confidencialidad & Anonimato del evaluador y protección de datos personales. \\
Trazabilidad & Cada caso registra responsable, fecha y bitácora de cambios. \\
Oportunidad & Notificaciones automáticas y control de plazos. \\
Eficiencia & Automatización y eliminación de intermediarios innecesarios. \\
Innovación & Uso de analítica y modelos predictivos para anticipar problemas. \\
Servicio & Foco en reducir la carga y mejorar la experiencia de los actores. \\
Mejora continua & Ciclos de retroalimentación a partir de las métricas. \\
\bottomrule
\end{tabularx}
\end{table}

\section{Objetivos Estratégicos}

\begin{enumerate}[leftmargin=1.4cm, itemsep=5pt]
  \item \textbf{OE1. Gestión de cupos eficiente.} Automatizar la solicitud,
        verificación de disponibilidad y asignación de cupos con criterios
        transparentes y tiempo de respuesta inferior a 24 horas.

  \item \textbf{OE2. Ciclo de vida de materias.} Formalizar el flujo de creación,
        edición, aprobación y publicación de materias con control de versiones.

  \item \textbf{OE3. Comunicación institucional unificada.} Proveer un canal
        único y trazable para la comunicación bidireccional entre estudiantes,
        docentes y secretaría.

  \item \textbf{OE4. Notificación temprana de faltas.} Registrar faltas en el
        momento y notificar al estudiante de forma automática e inmediata.

  \item \textbf{OE5. Evaluación docente oportuna.} Garantizar la aplicación,
        recordatorio y cierre de la evaluación docente dentro del periodo.

  \item \textbf{OE6. Gestión formal de quejas y sugerencias.} Categorizar,
        priorizar, asignar responsable y cerrar cada caso con trazabilidad total.

  \item \textbf{OE7. Fortalecimiento de la secretaría.} Dotar a la secretaría de
        herramientas de priorización, delegación y comunicación con los docentes.

  \item \textbf{OE8. Optimización de horarios.} Aplicar cuestionarios y modelos
        probabilísticos para programar horarios con mayor aceptación estudiantil.

  \item \textbf{OE9. Verificación institucional.} Ofrecer tableros de control a
        la secretaría y al departamento para verificar resultados, quejas y
        evaluaciones.

  \item \textbf{OE10. Analítica y soluciones rápidas.} Implementar un motor
        analítico que detecte anomalías y genere recomendaciones accionables.
\end{enumerate}

\section{Indicadores Clave de Desempeño}

\begin{table}[H]
\centering
\caption{Indicadores clave de desempeño (\ac{KPI})}
\label{tab:kpis}
\small
\begin{tabularx}{\textwidth}{@{}p{1.5cm} X p{1.7cm} p{2.0cm}@{}}
\toprule
\rowcolor{azulClaro}
\textbf{KPI} & \textbf{Definición} & \textbf{Meta} & \textbf{Frecuencia} \\
\midrule
TRC & Tiempo de respuesta a solicitud de cupo (horas) & $<24$ & Semanal \\
TRD & Tiempo de respuesta docente a mensaje (horas) & $<48$ & Semanal \\
TED & Tasa de evaluaciones docentes completadas & $100\%$ & Por periodo \\
TQR & Tasa de quejas resueltas dentro del plazo & $>90\%$ & Mensual \\
TMC & Tiempo medio de cierre de queja crítica (días) & $<3$ & Mensual \\
ICOM & Índice de satisfacción con la comunicación (1--5) & $>4.0$ & Por periodo \\
CNF & Cobertura de notificación de faltas el mismo día & $100\%$ & Semanal \\
ACE & Aceptación de horario asignado (1--5) & $>4.0$ & Por periodo \\
IDD & Índice de desempeño docente promedio (0--5) & $>4.0$ & Por periodo \\
EFA & Efectividad de acciones de la secretaría & $>85\%$ & Mensual \\
EAN & Exactitud del modelo de detección de anomalías & $>85\%$ & Trimestral \\
\bottomrule
\end{tabularx}
\end{table}

\section{Mapa Estratégico}

\begin{figure}[H]
\centering
\begin{tikzpicture}[
  cajita/.style={
    draw=azulOscuro, fill=azulClaro, thick, rounded corners=4pt,
    text=azulOscuro, align=center, font=\scriptsize,
    minimum width=7.5cm, minimum height=1.2cm, inner sep=4pt
  },
  meta/.style={
    draw=verdeUC, fill=verdeClaro, thick, rounded corners=4pt,
    text=verdeUC!60!black, align=center, font=\scriptsize,
    minimum width=7.5cm, minimum height=1.2cm, inner sep=4pt
  },
  node distance=0.5cm
]

\node[meta] (vision) {\textbf{VISIÓN}\\Ser referente universitario en gestión departamental};
\node[cajita, below=of vision] (mision) {\textbf{MISIÓN}\\Plataforma unificada, trazable y analítica};
\node[cajita, below=of mision] (perspectiva) {\textbf{PERSPECTIVA ESTRATÉGICA}\\Soluciones rápidas basadas en evidencia};
\node[cajita, below=of perspectiva] (resultados) {\textbf{RESULTADOS}\\KPI: TRC, TED, TQR, CNF, ACE, IDD};

\draw[flecha] (vision) -- (mision);
\draw[flecha] (mision) -- (perspectiva);
\draw[flecha] (perspectiva) -- (resultados);

\end{tikzpicture}
\caption{Mapa estratégico del \ac{SGDIC}: de la visión a los resultados medibles}
\label{fig:mapa-estrategico}
\end{figure}

\section{Beneficios Esperados por Actor}

\begin{table}[H]
\centering
\caption{Beneficios esperados por actor involucrado}
\label{tab:beneficios}
\small
\begin{tabularx}{\textwidth}{@{}l X@{}}
\toprule
\rowcolor{verdeClaro}
\textbf{Actor} & \textbf{Beneficios principales} \\
\midrule
Estudiante &
  Solicitud de cupos en línea con respuesta en menos de 24 h; notificación
  inmediata de faltas; canal formal y anónimo para evaluar y reportar; horarios
  más compatibles con su disponibilidad. \\
Docente &
  Gestión ágil de materias; canal único de mensajería; retroalimentación de
  estudiantes consolidada y anónima; menor carga administrativa. \\
Secretaría &
  Bandeja priorizada de solicitudes y quejas; asignación de responsables;
  comunicación formal con docentes; tableros de seguimiento; eliminación de
  registros duplicados en hojas de cálculo. \\
Departamento &
  Verificación consolidada de resultados, quejas y evaluaciones; indicadores en
  tiempo real; alertas tempranas; soporte objetivo para decisiones académicas. \\
Institución &
  Mayor satisfacción de la comunidad académica; trazabilidad auditable;
  experiencia replicable en otras unidades académicas. \\
\bottomrule
\end{tabularx}
\end{table}